\chapter{Abstract Interpretation}
\label{ch:ch11-abstract-interpretation}

Testing runs a program on some inputs, while optimization must reason about all inputs at once.
Enumerating every input is impossible, and sampling a few inputs proves nothing about the rest.
The problem of this chapter is to execute a program over descriptions of values rather than over values themselves, and to do so with a guarantee.
We replace concrete stores with abstract elements, concrete operations with sound counterparts, and concrete branching with joins over abstract paths.
When the abstract run says that no division by zero occurs, no execution from the described inputs divides by zero.
Section~\ref{sec:ch11-domains} builds the domains, Section~\ref{sec:ch11-soundness} states the guarantee, and Section~\ref{sec:ch11-worked} runs an interval analysis by hand on Figure~\ref{fig:ch11-branch-loop}.

\section{Concrete runs versus abstract runs}
\label{sec:ch11-domains}

A \term{concrete state} maps each variable to one exact value, for example a 32-bit integer or a pointer.
A \term{abstract state} maps each variable to a description that stands for many concrete values, for example an interval, a sign, or a parity class.
The \term{abstraction function} sends each concrete set to the smallest abstract element covering it.
The \term{concretization function} sends each abstract element back to the set of concrete states it describes.
They form a \term{Galois connection} when $\alpha(S)\sqsubseteq a$ if and only if $S\subseteq\gamma(a)$ for every concrete-state set $S$ and abstract element $a$.
Here $\alpha$ abstracts and $\gamma$ concretizes; the equivalence says that $\alpha(S)$ is the least abstract over-approximation of $S$.
We write the pair once per domain and reuse it for every program the domain analyzes.
The connection is the contract that lets an optimizer trust abstract results about concrete runs.

Consider tracking the sign of \code{x} through \code{x = y * z} with unknown inputs.
The concrete multiplication has four billion squared cases, far too many to enumerate.
For mathematical integers, abstract multiplication has a nine-entry table over negative, zero, and positive, extended with top and bottom for unknown and impossible.
Fixed-width arithmetic requires different rules when wraparound can change a sign.
If \code{y} is negative and \code{z} is positive, the abstract product is negative with no case split.
Multiplying top by zero yields zero; multiplying top by a nonzero sign yields top.
If either side is bottom, the product is bottom, which marks the path as unreachable rather than as mysterious.
The \enquote{execution} is an ordinary table lookup, yet its result covers all concrete executions it abstracts.
That coverage is the whole point, and the next section makes it a theorem shape rather than a slogan.

\begin{irlisting}
; Branch then loop: abstract states must cover both arms
; B1 reads n, B2 splits on n, B3 and B4 bound x, B5 loops
B1: n = read_input()
B1: br n < 0, B2, B3
B2: x = 0 - n
B2: jmp B4
B3: x = n
B3: jmp B4
B4: i = 0
B4: jmp B5
B5: br i < x, B6, B7
B6: i = i + 1
B6: jmp B5
B7: use i
\end{irlisting}

The listing motivates intervals over signs, because the loop bound \code{x} is nonnegative on both arms but neither arm alone proves it.
A sign domain reports \code{x} as nonnegative only if it has a dedicated element for that property.
For mathematical integers, the negative arm yields $[1,+\infty)$ after negation, while the other arm yields $[0,+\infty)$; their interval join is $[0,+\infty)$.
The join at $B4$ therefore keeps exactly the fact the loop needs: the bound cannot be negative.
A parity domain would instead report nothing useful here, because nonnegativity is invisible to it.
No single domain serves every client, which is why this chapter treats domain choice as an engineering decision with a soundness proof obligation.
Algorithm~\ref{alg:ch11-abstract} runs any domain that meets the obligation.

\section{Soundness and what it does not promise}
\label{sec:ch11-soundness}

An abstract transfer is \term{sound} when every concrete behavior it abstracts stays inside the abstract result.
Formally, abstracting the concrete outputs from any covered input set must sit below the abstract transfer applied to the abstraction of that input set.
Soundness composes across blocks: sound steps chained by joins stay sound end to end.
The fixed point computed by Algorithm~\ref{alg:ch11-abstract} therefore over-approximates every reachable concrete state at every block.
An optimizer may delete a check that the over-approximation proves redundant, because no omitted concrete state could have needed it.
It may not insert a behavior that only the approximation exhibits, because the approximation can contain states no execution reaches.
An over-approximation can prove universal properties, such as every returned value being nonnegative; an abstract witness alone cannot prove that a concrete execution exists.

Soundness does not imply precision, completeness, or speed.
An analysis that reports top everywhere is sound and useless, since it rules out nothing.
An analysis that tracks exact values along every path is precise but as expensive as testing every input.
An analysis over a tall domain may need widening to terminate, trading precision for a bound both analyses share.
Clients must therefore read abstract results as permissions to act when the approximation is strong, not as obligations to act whenever it is weak.
A division guard stays when intervals cannot exclude zero, even when testing never divides by zero.
A bounds check stays when the index interval escapes the array extent by even one abstract value.
The compiler reports \enquote{unknown} by keeping the check, which is the safe default that soundness demands.
A frequent misunderstanding equates bottom with failure of the analysis.
Bottom at a block means that no concrete execution described by the entry abstraction reaches that block with those abstract inputs.
Dead branches, contradictory guards, and empty input descriptions all produce bottom legitimately.
An optimizer should feed bottom forward as unreachability, not restart the analysis hoping for a friendlier answer.
A second misunderstanding treats a sound interval as a guarantee that every value inside occurs.
The interval $[0, 10]$ for a loop counter does not promise that $7$ occurs, only that nothing outside $[0, 10]$ occurs.
Loop peeling based on the second reading introduces behaviors the program never had.
The safe readings all flow from the Galois connection in one direction: concrete sets inside abstract descriptions, never the reverse.

\begin{algorithm}
\caption{Abstract interpretation over a CFG with widening at loop heads}\label{alg:ch11-abstract}
\begin{pseudocode}
\Require reachable CFG with predecessor-free entry, monotone sound edge transfers $F_{P,B}$, widening $\nabla$, widening points meeting every cycle
\Ensure over-approximation $X[B]$ of all reachable concrete states at each $B$
\State $X[B] \gets \bot$ for every block $B$
\State $X[entry] \gets$ abstraction of the known input description
\State $W \gets$ all blocks in $G$
\While{$W \neq \emptyset$}
\State remove some $B$ from $W$
\State $in \gets \join_{P \in \Pred(B)} F_{P,B}(X[P])$, or entry abstraction at entry
\State $new \gets X[B] \join in$
\If{$B$ is a widening point and $new \not\sqsubseteq X[B]$}
\State $new \gets X[B] \nabla new$
\EndIf
\If{$new \neq X[B]$}
\State $X[B] \gets new$
\State add $\Succ(B)$ to $W$
\EndIf
\EndWhile
\end{pseudocode}
\end{algorithm}

Algorithm~\ref{alg:ch11-abstract} is the lattice solver from the previous chapter with two abstract-interpretation additions.
First, the seed is the abstraction of the caller-supplied input description rather than a generic extreme.
Second, widening points apply widening when the new value climbs outside the stored value.
Each edge transfer executes the predecessor block and refines by the condition selecting that edge; without this refinement the worked branch example would lose its bounds.
Natural-loop headers suffice for reducible CFGs; an irreducible graph needs widening points that also meet its other cycles.
Every other block uses the plain join, so precision is lost only where loops demand it.
Termination follows because widening chains are finite by construction and non-head blocks move along finite joins between widenings.
Soundness follows because seeds cover the inputs, transfers are sound, joins cover both arms, and widening only climbs to coarser covers.
The result $X[B]$ therefore contains every concrete state that any described execution can carry into $B$.
Section~\ref{sec:ch11-worked} checks this claim on a concrete trace below.

Narrowing can recover precision after widening overshoots, at the cost of a second pass.
Let $F$ be the global monotone equation operator, including the entry boundary, and let $X$ be the converged widened result with $F(X)\sqsubseteq X$.
A synchronous replay $X'=F(X)$ remains an over-approximation of the least fixed point and satisfies $F(X')\sqsubseteq X'$ by monotonicity.
An arbitrary intersection with a smaller state is not justified merely because the old state was sound.
Descending chains need not be finite, so a fixed replay budget or a narrowing operator with a termination argument is required.
An implementation should budget narrowing explicitly, for example one descent per loop head.
Unbounded narrowing reintroduces the termination problem that widening solved.
Morgan discusses how production optimizers budget such refinement passes against compile-time limits, see~\cite{morgan-1998-optimizing}.
\section{Worked example with intervals}
\label{sec:ch11-worked}

We run intervals by hand over Figure~\ref{fig:ch11-branch-loop} with input $n$ in $[-5, 5]$.
The domain tracks one interval per variable, with top as unbounded and bottom as empty.
Transfer for \code{x = 0 - n} negates the interval, transfer for \code{x = n} copies it, and transfer for \code{i = i + 1} shifts both bounds up by one.
The branch on \code{n} splits the input interval into $[-5, -1]$ on the true arm and $[0, 5]$ on the false arm.
Negation maps the true arm to $[1, 5]$, while the false arm keeps $[0, 5]$.
The join at $B4$ is $[0, 5]$, which proves the loop bound nonnegative on every path.
Table~\ref{tab:ch11-intervals} records the abstract states that justify each step.

\begin{figure}
\begin{tikzpicture}
\node[basic block] (b1) {B1: n = input};
\node[basic block, below left=1.1cm and 0.7cm of b1] (b2) {B2: x = 0 - n};
\node[basic block, below right=1.1cm and 0.7cm of b1] (b3) {B3: x = n};
\node[basic block, below=2.4cm of b1] (b4) {B4: i = 0};
\node[basic block, below=1.1cm of b4] (b5) {B5: test i less x};
\node[basic block, below left=1.1cm and 0.7cm of b5] (b6) {B6: i = i + 1};
\node[basic block, below right=1.1cm and 0.7cm of b5] (b7) {B7: use i};
\draw[cfg edge] (b1) -- (b2);
\draw[cfg edge] (b1) -- (b3);
\draw[cfg edge] (b2) -- (b4);
\draw[cfg edge] (b3) -- (b4);
\draw[cfg edge] (b4) -- (b5);
\draw[cfg edge] (b5) -- (b6);
\draw[cfg edge] (b5) -- (b7);
\draw[back edge] (b6) -- (b5);
\end{tikzpicture}
\caption{Branch followed by a loop, where the join at $B4$ must recover nonnegativity of \code{x} before the loop head at $B5$ analyzes the bound.}
\label{fig:ch11-branch-loop}
\end{figure}

\begin{table}
\caption{Interval facts for the branch-join prefix with $n$ in $[-5, 5]$, showing how negation and join recover a nonnegative bound.}
\label{tab:ch11-intervals}
\begin{tabular}{lcc}
\toprule
Block & Abstract entry state & Abstract exit state \\
\midrule
B1 & n in [-5, 5] & n in [-5, 5] \\
B2 & n in [-5, -1] & x in [1, 5] \\
B3 & n in [0, 5] & x in [0, 5] \\
B4 & x in [0, 5] & i in [0, 0], x in [0, 5] \\
\bottomrule
\end{tabular}
\end{table}

The loop head $B5$ starts with \code{i} in $[0, 0]$ and \code{x} in $[0, 5]$.
The first body visit shifts \code{i} to $[1, 1]$, and the head join becomes $[0, 1]$.
Without widening the next visits produce $[0, 2]$, $[0, 3]$, and so on up to $[0, 5]$, each trip adding one abstract value.
With widening after two climbs, the head jumps from $[0, 1]$ to $[0, +\infty)$ on the third visit.
Recomputing the loop equation from the widened invariant refines the true edge by \code{i < x}, giving $i\in[0,4]$ at body entry and $i\in[1,5]$ after increment.
Joining that back-edge result with the initial $[0,0]$ recovers $[0,5]$ at the header.
This is one replay of the loop equation; a block-by-block worklist may need several removals to carry the refined result around the loop.
The final abstract state at $B7$ has \code{i} in $[0, 5]$, which covers every concrete trip count for every input in $[-5, 5]$.
The concrete body executes at most five times; the six values $0$ through $5$ include the final header visit that exits.
Neither a range of counter values nor an upper bound on trips permits unconditionally executing that many body copies.

A concrete check confirms coverage on two representative inputs.
With $n = -3$ the program takes the true arm, sets \code{x} to $3$, and iterates \code{i} through $0$, $1$, $2$, $3$.
With $n = 4$ it takes the false arm, sets \code{x} to $4$, and iterates through $0$ to $4$.
Both traces stay inside the abstract states of Table~\ref{tab:ch11-intervals} at every block.
No trace exhibits \code{x} negative, \code{i} above $5$, or entry to $B6$ with \code{i} outside $[0, 5]$.
The abstract run proves these absences for the remaining inputs without executing them.
That is the return on the extra domain machinery: one run covers the whole input interval.

\section{Implementation notes and cost}
\label{sec:ch11-implementation}

An interval domain stores two bounds per variable with flags for open and unbounded ends.
Join takes componentwise minima of low bounds and maxima of high bounds, each in constant time.
Meet takes maxima of lows with minima of highs and yields bottom when low exceeds high.
Transfer for addition shifts both ends, transfer for negation swaps and negates them, and transfer for multiplication takes minima and maxima over the four products of ends.
Guards refine states by intersecting with the guard condition, for example \code{i} below \code{x} caps the high end of \code{i}.
Memory cost is two bounds per variable per block, and time per visit is linear in the number of live variables times constant domain operations.
Appel shows how such domains plug into the same worklist machinery as simpler flow facts, see~\cite{appel-1998-modern}.

Let $V$ be the variable count and $N$ the block count.
If block $B$ is visited $r_B$ times, scanning its incoming edges costs $O(r_B|\Pred(B)|V)$ interval operations, in addition to executing the edge transfers.
Sum these costs over blocks and add the cost of narrowing replays.
The widening delay alone does not bound $r_B$: nested cycles, queue order, and propagation between widening points also matter.
Memory is $O(NV)$ intervals, excluding the CFG, with constant-size bounds assumed.
Precision depends on where guards refine and where widening fires, not only on domain height.
Relational domains that track facts like \code{i} below \code{x} keep bounds the non-relational domain loses, at quadratic or worse cost.
The C standard leaves signed overflow, pointer provenance, and trap behavior to a mix of defined, undefined, and implementation choices, so the abstract transfer must encode those rules rather than ideal arithmetic, see~\cite{wg14-2011-c11}.
For the primary development of abstraction, fixed points, and sound approximation, see Cousot and Cousot~\cite{cousot-1977-abstract}.

\section{Exercises}
\label{sec:ch11-exercises}

\begin{enumerate}
\item \label{ex:ch11-sign-table} Write the full abstract multiplication table for the sign domain with top and bottom.
Include all nine sign pairs plus the rows and columns for top and bottom, and justify each entry by concretization coverage.
Then exhibit one program where sign analysis proves absence of division by zero and one where it must keep the guard.
Run both on Figure~\ref{fig:ch11-branch-loop} adapted to divide by \code{x} instead of looping on it.
\item \label{ex:ch11-unsound-twin} Take Algorithm~\ref{alg:ch11-abstract} and replace the join at $B4$ with an intersection, keeping everything else identical.
Show on Table~\ref{tab:ch11-intervals} that the result claims \code{x} in $[1, 5]$ and drops the concrete behavior with $n = 0$.
Explain in safety terms why the mutated analysis is unsound and which deleted guard could then miscompile.
Propose a one-line fix and the test that would have caught the mutation.
\item \label{ex:ch11-widen-budget} Implement intervals for Figure~\ref{fig:ch11-branch-loop} with configurable widening delay $W$.
Measure head visits and final bound width for $W$ set to $1$, $2$, $5$, and unbounded, with input $n$ in $[-5, 5]$.
Add one narrowing replay and report how much width each $W$ recovers.
Recommend a default $W$ with reasons covering time, precision, and predictability.
\end{enumerate}

\section{Further reading}
\label{sec:ch11-reading}

Appel presents program analysis as executable approximation with concrete transfer code, see~\cite{appel-1998-modern}.
Morgan explains how analyses of this form feed loop and redundancy optimizations, see~\cite{morgan-1998-optimizing}.
Safonov examines what a soundness argument must cover before a compiler may act on it, see~\cite{safonov-2010-trustworthy}.
The C language definition determines which concrete behaviors the abstraction must cover, see~\cite{wg14-2011-c11}.
